%% kapa: exam=1 round=31 year=2020 period=1 form=A subject=real_estate subject_ko=부동산학원론 seq=6 number=86
%% source: Q-Net 감정평가사 기출문제 (cst00309) artlSeq=5207219 file=extracted/1008/5207219/A형/A형/2020년도 제31회 감정평가사 1차 시험_1교시_A형.pdf
%% conversion: kordoc 4.24.51 -> tools/gen1.py (auto); review=proofread-sonnet; flags=none
\begin{question}{86}
감정평가에 관한 규칙상 용어의 정의로 옳지 않은 것은?
\begin{choices}
  \item 기준시점이란 대상물건의 감정평가액을 결정하는 기준이 되는 날짜를 말한다.
  \item 가치형성요인이란 대상물건의 경제적 가치에 영향을 미치는 일반요인, 지역요인 및 개별요인 등을 말한다.
  \item 동일수급권이란 대상부동산과 대체ㆍ경쟁 관계가 성립하고 가치 형성에 서로 영향을 미치는 관계에 있는 다른 부동산이 존재하는 권역을 말하며, 인근지역과 유사지역을 포함한다.
  \item 임대사례비교법이란 대상물건과 가치형성요인이 같거나 비슷한 물건의 임대사례와 비교하여 대상물건의 현황에 맞게 사정보정, 시점수정, 가치형성요인 비교 등의 과정을 거쳐 대상물건의 임대료를 산정하는 감정평가방법을 말한다.
  \item 수익분석법이란 대상물건이 장래 산출할 것으로 기대되는 순수익이나 미래의 현금흐름을 환원하거나 할인하여 대상물건의 가액을 산정하는 감정평가방법을 말한다.
\end{choices}
\end{question}
\answer{5}
